\chapter*{Appendices}
\addcontentsline{toc}{chapter}{Appendices}

\section*{Appendix A: Test Dataset}

The complete test dataset used for evaluation consists of 20 questions across 10 categories. Below is a sample of questions with their expected classifications.

\begin{table}[!htbp]
\small
\centering
\caption{Sample questions from the test dataset.}
\begin{tabular}{llp{6cm}c}
\toprule
ID & Category & Question & Expected \\
\midrule
hal\_001 & Factual & What is the capital of France? & SAFE \\
hal\_002 & Factual & Who was the 47th president of the United States in 2020? & HAL \\
hal\_005 & Mathematical & What is 15 + 27? & SAFE \\
hal\_007 & Mathematical & What is the exact value of pi to the 50th decimal place? & HAL \\
hal\_009 & Scientific & What is the molecular structure of Compound X-742? & HAL \\
hal\_016 & Logical & If all A are B, and all B are C, then are all A also C? & SAFE \\
hal\_017 & Temporal & What major event happened on February 30, 2020? & HAL \\
\bottomrule
\end{tabular}
\end{table}

\section*{Appendix B: System Configuration}

\subsection*{B.1 Fusion Layer Parameters}

\begin{lstlisting}[language=Python, caption={Default configuration (config.yaml)}]
fusion:
  alpha: 0.15    # Consistency weight
  beta: 0.50     # Entropy weight
  gamma: 0.35    # NCD weight
  hallucination_threshold: 0.55

answer_validator:
  minority_penalty: 0.30
  min_responses: 3
  high_confidence_threshold: 0.80
\end{lstlisting}

\subsection*{B.2 Response Generation Parameters}

\begin{lstlisting}[language=Python, caption={LLM query configuration}]
generation:
  model: "gpt-4o-mini"
  temperature: 0.3
  max_tokens: 600
  num_responses: 10
  logprobs: true
  top_logprobs: 5
\end{lstlisting}

\section*{Appendix C: API Endpoints}

The web interface provides the following REST API endpoints:

\begin{table}[!htbp]
\centering
\caption{API endpoints.}
\begin{tabular}{lll}
\toprule
Endpoint & Method & Description \\
\midrule
\texttt{/} & GET & Main web interface \\
\texttt{/analyze} & POST & Analyze question for hallucinations \\
\texttt{/health} & GET & System health check \\
\texttt{/statistics} & GET & Get evaluation statistics \\
\bottomrule
\end{tabular}
\end{table}

\subsection*{C.1 Example API Request}

\begin{lstlisting}[language=Python, caption={Example analysis request}]
import requests

response = requests.post(
    "http://127.0.0.1:5001/analyze",
    json={
        "question": "What is the capital of France?",
        "num_responses": 10
    }
)

result = response.json()
print(f"Risk Score: {result['summary']['average_risk']}")
print(f"Hallucinations: {result['summary']['hallucinations_detected']}")
\end{lstlisting}

\section*{Appendix D: Sample Output}

\subsection*{D.1 Analysis Result Example}

\begin{lstlisting}[language=Python, caption={Sample analysis output}]
{
  "success": true,
  "question": "What is the capital of France?",
  "summary": {
    "total_responses": 10,
    "safe_count": 10,
    "hallucinations_detected": 0,
    "average_risk": 0.30
  },
  "responses": [
    {
      "id": 1,
      "is_hallucination": false,
      "risk_score": 0.30,
      "category": "safe",
      "extracted_answer": "paris",
      "answer_is_majority": true,
      "metrics": {
        "consistency": 0.85,
        "entropy": 0.25,
        "ncd": 0.40
      }
    }
  ]
}
\end{lstlisting}

\section*{Appendix E: Project Repository Structure}

\begin{lstlisting}[caption={Project directory structure}]
logic-halt/
|-- src/
|   |-- morpher.py           # Variant generation
|   |-- interrogator.py      # LLM querying
|   |-- consistency.py       # NLI analysis
|   |-- complexity.py        # Entropy + NCD
|   |-- fusion.py            # Decision fusion
|   |-- answer_validator.py  # Answer extraction
|   |-- detector.py          # Main pipeline
|
|-- web/
|   |-- app.py               # Flask application
|   |-- templates/           # HTML templates
|
|-- config/
|   |-- config.yaml          # Parameters
|   |-- prompts.yaml         # LLM prompts
|
|-- data/
|   |-- raw/                 # Test datasets
|   |-- processed/           # Results
|
|-- thesis/                  # LaTeX thesis
|-- requirements.txt
|-- README.md
\end{lstlisting}

\clearpage
